% Vietnamese translation for OI Public Library
% Source-File: IOI中国国家候选队论文/2009/汤可因/浅析竞赛中一类数学期望问题的解决方法.ppt
\documentclass[11pt]{article}
\usepackage{oipl-vn}

\title{Bàn sơ về phương pháp giải một lớp bài toán kỳ vọng trong thi đấu\\{\large Ghi chú theo slide}}
\author{Tang Keyin\\Trường Trung học số 8 Phúc Châu, Phúc Kiến}
\date{2009}

\begin{document}
\maketitle

\origtitle{Methods for expectation problems in contests}
\sourcepath{IOI-China-Candidate-Team-Papers/2009/Tang-Keyin/expectation-problems-slides.ppt}

\section{Kiến thức chuẩn bị}

Slide nhấn mạnh ba công cụ: tính tuyến tính của kỳ vọng, xác suất toàn phần và
kỳ vọng có điều kiện. Trong nhiều bài, ta không cần biết toàn bộ phân phối của
biến ngẫu nhiên; chỉ cần viết đúng quan hệ kỳ vọng giữa các trạng thái.

\section{Hai phương pháp chính}

\begin{enumerate}
  \item Truy hồi hoặc quy hoạch động: trạng thái có hướng đi rõ ràng, nên kỳ
  vọng của trạng thái hiện tại được tính từ các trạng thái nhỏ hơn hoặc gần
  hơn đích.
  \item Lập hệ phương trình tuyến tính: các trạng thái có thể quay lại nhau,
  tạo chu trình phụ thuộc, nên cần viết một phương trình cho mỗi trạng thái
  rồi giải hệ.
\end{enumerate}

\section{Ví dụ}

\textit{Congcong và Coco}, \textit{Highlander} và \textit{RedIsGood} thuộc
nhóm truy hồi hoặc quy hoạch động. \textit{First Knight} và \textit{Mario}
minh họa trường hợp phụ thuộc vòng, nơi phương pháp hệ tuyến tính tự nhiên
hơn.

\section{Ghi chú}

Khi đọc slide, hãy xác định trạng thái dừng trước, sau đó kiểm tra tổng xác
suất chuyển ra khỏi mỗi trạng thái. Việc viết phương trình đúng thường quan
trọng hơn việc chọn kỹ thuật giải hệ.

\end{document}
